\input{../tex-inputs/latex-leseansicht-vorspann}

\section[Arthur Schnitzler an Gustav Schwarzkopf, {[}24. 4. 1908?{]}]{L04199 Arthur Schnitzler an Gustav Schwarzkopf, {[}24. 4. 1908?{]}}
\nopagebreak\mylabel{L04199v}
\rehead{ }\normalsize\beginnumbering\briefempfaengerindex{Schwarzkopf, Gustav@\textsc{Schwarzkopf, Gustav}!zzzSchnitzler, Arthur@\emph{von Arthur Schnitzler}!1908-04-241@{{[}24. 4. 1908?{]}}|(be}
\toendnotes[C]{\smallbreak\pagebreak[2]}
\correspDesc{Versand  durch Arthur Schnitzler am [24. 4. 1908?] in Edmund-Weiß-Gasse
\newline{}Erhalt  durch Gustav Schwarzkopf am [24. 4. 1908?] in Tiefer Graben 23}\toendnotes[C]{\smallbreak}
\Standort{CUL, Schnitzler, B 96.}
\physDesc{Visitenkarte, 125 Zeichen
\newline{}Handschrift: Bleistift, deutsche Kurrent}\toendnotes[C]{\smallbreak}
\pstart
           \noindent{}{\pb}Wir reiſen vorläufig nicht.\pend
           
\pstart
           \centering{}\textcolor{gray}{\textbf{Dr Arthur Schnitzler}}\pend
           
\pstart
           \raggedleft{}\textcolor{gray}{\textbf{Wien\oindex{Wien@\textbf{Wien}|pw}}}\pend
           
\pstart
           Wenn Sie \label{K_L04199-1v}\edtext{heute Freitag}{\lemma{\textnormal{\emph{heute Freitag}}}\Cendnote{\textnormal{Die Visitenkarte ist undatiert. Eine weitere Visitenkarte dieses Typs (bei der
                  die beiden Unterstriche unter der Hochstellung bei »D\textsuperscript{r}« waagrecht verlaufen) verwendete Schnitzler
                  am XXXX Auszeichnungsfehler: Dokument L03809 nicht gefunden, so dass eine zeitliche Einordnung möglich ist.
                  Gleicht man das mit dem \emph{Tagebuch}\pwindex{Schnitzler, Arthur 15. 5. 1862 Wien – 21. 10. 1931 ebd.@\textsc{Schnitzler, Arthur} (15. 5. 1862 Wien – 21. 10. 1931 ebd.)!Tagebuch@\strich\emph{Tagebuch}|pwk} ab, so dürfte die
                  vorliegende Karte am 24. 4. 1908 verfasst
                  sein, da an diesem Tag alle in der Frankgasse 1\oindex{Wien@\textbf{Wien}!IX., Alsergrund@\textbf{IX., Alsergrund}!Frankgasse 1@\textbf{Frankgasse 1}|pwk} zusammentrafen. Am nächsten
                  Tag reisten Arthur und Olga Schnitzler\pwindex{Schnitzler, Olga 17.\,1.\,1882 Wien – 13.\,1.\,1970 Lugano@\textsc{Schnitzler, Olga} (17.\,1.\,1882 Wien – 13.\,1.\,1970 Lugano)|pwk}
                  ab.}}}\label{K_L04199-1} abend
            Frankgaſſe\oindex{Wien@\textbf{Wien}!IX., Alsergrund@\textbf{IX., Alsergrund}!Frankgasse 1@\textbf{Frankgasse 1}|pw} kämen, würde ſich
      (abgeſehen von den andern) Mama\pwindex{Schnitzler, Louise 8.\,7.\,1840 Kőszeg – 9.\,9.\,1911 Wien@\textsc{Schnitzler, Louise} (8.\,7.\,1840 Kőszeg – 9.\,9.\,1911 Wien)|pwv} ſehr
      freuen. –\pend
           \selectlanguage{ngerman}\endnumbering\briefempfaengerindex{Schwarzkopf, Gustav@\textsc{Schwarzkopf, Gustav}!zzzSchnitzler, Arthur@\emph{von Arthur Schnitzler}!1908-04-241@{{[}24. 4. 1908?{]}}|)be}\mylabel{L04199h}
\begin{anhang}
\end{anhang}\newcommand{\dateiname}{L04199}\newcommand{\titel}{Arthur Schnitzler an Gustav Schwarzkopf, [24. 4. 1908?]}\newcommand{\editorInnen}{Herausgegeben von Jahnke, SelmaMüller, Martin Anton}\input{../tex-inputs/latex-leseansicht-abspann}
